\chapter{TEKNOLOGI DAN ASUMSI}

\section{Teknologi Terpilih}
Implementasi menggunakan Go dengan modul minimum Go 1.23. Driver PostgreSQL adalah \texttt{pgx/v5}, broker diakses melalui \texttt{nats.go}, dan validasi fixture kontrak memakai JSON Schema Draft 2020-12. Compose memilih PostgreSQL 16 dan NATS 2.10; tag tersebut belum mengunci image digest. Auth menggunakan library JWT dan SQLite, sedangkan masing-masing consumer memiliki SQLite privat. Pilihan ini memisahkan canonical data dari state autentikasi dan idempotensi consumer.

Go dipakai untuk HTTP, parsing JSON, goroutine, dan pembatalan melalui \texttt{context}. PostgreSQL menyediakan transaksi untuk perubahan hazard, riwayat, dan outbox; JSONB menjaga atribut dinamis tanpa mengubah kolom kanonik [3,4,5]. JetStream menyediakan persistent stream dan durable consumer, bukan jaminan exactly-once menyeluruh [6,7]. Docker Compose mengelola service terpisah pada satu host [9].

\section{Alternatif Teknologi}
\begin{itemize}
\item \textbf{PostgreSQL + JSONB}: kolom bertipe menjaga identitas dan constraint, sedangkan atribut aditif tetap berupa dokumen. Dibanding seluruh atribut menjadi kolom, tambahan field tidak membutuhkan migrasi; dibanding document store terpisah, transaksi latest/history/outbox tetap berada pada satu store.
\item \textbf{JetStream}: menambahkan persistence, replay, dan posisi durable pada NATS. NATS Core tidak menyimpan backlog untuk consumer offline. Kafka atau RabbitMQ juga dapat mendukung distribusi event, tetapi menambah pilihan konfigurasi yang belum diperlukan untuk demonstrasi satu host ini.
\item \textbf{SQLite privat}: sesi refresh dan inbox tetap tersimpan saat restart, tanpa server database tambahan. Penyimpanan memory lebih sederhana tetapi kehilangan state; beberapa replica yang harus berbagi state memerlukan rancangan database dan koordinasi berbeda.
\item \textbf{Docker Compose}: cukup untuk container, network, dan volume pada satu host. Kubernetes belum diperlukan untuk cakupan demo; konfigurasi Compose ini juga tidak memberikan failover lintas mesin.
\end{itemize}

\section{Asumsi dan Batas Implementasi}
\begin{enumerate}
\item Mock menggunakan fixture sintetis, bukan data produksi BMKG/PVMBG. Nama dan koordinat gunung pada referensi Aggregator diberi label simulasi Tamriel; tidak mewakili lokasi gunung sebenarnya.
\item Dua loop polling independen berjalan segera saat startup, lalu mengikuti interval default 5 detik. Satu sumber hanya memiliki satu siklus aktif; tick saat sibuk dilewati. BMKG mengambil seismic dan warning secara concurrent.
\item Baseline ingest menggunakan full snapshot, bukan cursor \texttt{since}. Endpoint mock tetap mendukung \texttt{since}; full snapshot dipilih agar warning terlambat dan perubahan field pada ID lama tidak terlewat. Biaya baca bertambah bersama riwayat sumber.
\item Hanya perubahan aditif yang didukung otomatis. Field wajib hilang, tipe salah, enum tak dikenal, duplicate JSON key, atau volcano ID tak dikenal ditolak per record; record valid tetap diproses.
\item Baseline memakai satu Aggregator, satu database, dan satu node broker. Replikasi, failover lintas host, TLS produksi, dan exactly-once global berada di luar klaim implementasi ini.
\end{enumerate}

\section{Konfigurasi dan Observabilitas}
Repository menyediakan \texttt{.env.example} tanpa secret aktif; credential lokal tidak masuk Git. Setiap service menerima subset konfigurasi miliknya. API key BMKG, token PVMBG, internal key Aggregator, dan token NATS dipisahkan.

\begin{center}
\small
\begin{tabularx}{\textwidth}{>{\raggedright\arraybackslash}p{0.47\textwidth} >{\raggedright\arraybackslash}X}
\toprule
Parameter Aggregator & Default implementasi \\
\midrule
Interval poll / stale threshold & 5 s / 15 s \\
Timeout HTTP BMKG / PVMBG & 1 s / 4 s \\
Batas respons sumber / envelope event & 1 MiB / 64 KiB \\
Pool baca / tulis BMKG / tulis PVMBG / relay & 10 / 2 / 2 / 2 koneksi \\
Deadline query baca / transaksi tulis & 200 ms / 2 s \\
Admission API internal & 128 request data aktif \\
Relay poll / batch / pending maksimum & 500 ms / 20 pesan / 10.000 pesan \\
Timeout publish / rentang retry & 2 s / 1--30 s dengan jitter \\
\bottomrule
\end{tabularx}
\end{center}
Nilai tersebut merupakan konfigurasi awal, bukan hasil benchmark. \texttt{/health} menunjukkan liveness; \texttt{/ready} menilai dependency sesuai peran service. Kegagalan sumber atau broker tidak dengan sendirinya membuat API baca Aggregator unready selama database masih dapat dibaca. Log terstruktur menyertakan operasi, correlation ID, durasi, dan status tanpa payload mentah atau credential. Saat shutdown, context dibatalkan dan worker poller serta relay ditunggu sebelum resource ditutup.

\section{Kontrak Mock dan Data Awal}
BMKG melayani \texttt{/seismic-events} dan \texttt{/tsunami-warnings} dengan header \texttt{X-BMKG-Key}; delay default 100 ms. PVMBG melayani \texttt{/volcanic-reports} dengan Bearer token dan delay acak 500--3.000 ms. Endpoint admin \texttt{/admin/schema-version}, \texttt{/admin/outage}, serta tambahan \texttt{/admin/delay} menggunakan credential PVMBG yang sama. Credential diperiksa sebelum delay.

Kedua mock memulai dengan 20 record historis dan menambah record setiap 10 detik pada konfigurasi default. ID memuat run ID agar restart mock tidak menimpa riwayat lama. Toggle PVMBG v2 menambahkan \texttt{confidence\_level} pada serialisasi seluruh laporan; outage memberi 503 untuk data, tetapi endpoint admin dan liveness tetap dapat diakses.

\section{Menjalankan dan Memverifikasi}
Definisi Compose di repository sudah mencakup sembilan service utama; consumer ketiga tersedia melalui profil demo. Setup mencakup secret berbeda dan key JWT sebelum menjalankan \texttt{docker compose up --build}. Pengujian deployment penuh dan tes komponen native memakai lingkungan berbeda.

Komponen diverifikasi dengan tes unit/HTTP, race detector, build, vet, smoke kedua mock, serta tes Store dan publisher terhadap PostgreSQL/NATS disposable. Pemeriksaan terbaru pada commit \texttt{19beead} lulus. Dependency native yang dipakai adalah PostgreSQL 18.6 dan NATS 2.10.24; karena Compose menargetkan PostgreSQL 16, hasil ini bukan pengganti verifikasi image Compose. Rincian mapping dan storage ada pada Bab 4 dan 7. Evidence stack Compose dan pengujian lintas layanan tersedia pada lingkungan berbeda [10]. Bukti tersebut dipisahkan dari tes native; keberadaan evidence tidak berarti seluruh acceptance telah diulang pada commit main atau setiap profil volume data.

\subsection{Urutan Menjalankan Demo}
Dari root repository, salin \texttt{.env.example} menjadi \texttt{.env}, isi credential yang berbeda untuk setiap domain dan client, serta tentukan \texttt{COMPOSE\_PROJECT\_NAME}. Parameter biasa dapat memakai default Compose. Buat pasangan RSA pada \texttt{.secrets/}, lalu muat ke volume private/public yang terpisah. Langkah lengkap berada di \texttt{scripts/README.md}; jangan mengganti key yang sudah dipakai tanpa mempertimbangkan token aktif.
\begin{quote}\small\ttfamily
python3 scripts/setup-jwt-keys.py\\
docker compose config --quiet\\
docker compose up --build -d --wait
\end{quote}
Healthcheck menilai kesiapan startup, bukan seluruh acceptance. Login dilakukan ke Auth dengan identitas client yang sesuai sebelum mengakses API. Untuk menambahkan consumer ketiga tanpa membangun ulang producer:
\begin{quote}\small\ttfamily
docker compose --profile demo up -d --no-deps third-consumer
\end{quote}
Pengujian integrasi dijalankan melalui \texttt{scripts/verify-integration.py} dengan opsi \texttt{--scenario all}. Script tersebut mengubah schema/outage, menghentikan service sementara, dan membangun ulang consumer; gunakan stack khusus eksperimen. Demo console dapat membantu operator menjalankan command, tetapi hanya tersedia di branch tooling demo dan tidak diperlukan untuk jalur CLI.

\subsection{Lingkungan Evidence}
Evidence mencatat trial 5 Oktober 2026 pada commit \texttt{071b32d}, yaitu branch tooling demo, bukan HEAD main \texttt{19beead} [10]. Lingkungan trial adalah WSL2 Linux, 12 CPU logical, Docker 27.0.3, Compose 2.28.1, PostgreSQL 16.14, dan NATS 2.10.29. Build image aplikasi memakai Go 1.23; pemeriksaan lokal menggunakan Go 1.26.4. Port host API sampai Auth memakai 8280--8284. Seed 20 record, interval event 10 detik, polling 5 detik, serta deadline write 2 detik tetap merupakan profil awal, bukan jaminan skalabilitas.

